% ============================================================================
% Fichier  : partie2/P02_C04_ebios_risk_manager.tex
% Titre    : EBIOS Risk Manager
% Auteur   : MINKA MI NGUIDJOÏ Thierry Emmanuel
% Version  : 1.0 -- novembre 2026
% Corrigés : corriges/Corr_C04.tex
% ============================================================================
% Vérifié sur cyber.gouv.fr (consulté le 05/11/2026) : noms des cinq
% ateliers ; synthèse « conformité » et « scénarios » ; socle de sécurité ;
% scénarios centrés sur les menaces intentionnelles et ciblées ; guide
% publié le 9 octobre 2018 avec le Club EBIOS. Version 1.5 de mars 2024
% signalée par Wikipédia (lien vers le guide) : à confirmer sur le guide.
% NON RELU dans le guide (PDF non téléchargé) et signalé à l'auteur : la
% formule du niveau de menace d'une partie prenante (dépendance ×
% pénétration / maturité × confiance), la séquence « connaître, rentrer,
% trouver, exploiter », le nom « plan d'amélioration continue de la
% sécurité », les deux cycles. Les échelles de gravité et de vraisemblance
% du guide ne sont pas citées : le chapitre emploie l'échelle du cours.
% Faits du cas (bible) : LiZbiZ, quatre centres de données (Douala,
% Yaoundé, Garoua, Bafoussam), douze agences, application intégrée ;
% M. ETEME, Mme NKENG, M. KOUNOU, M. BEYALA (AUDIX CONSEIL) ; TRANSLOG
% (112,7 M FCFA, OWONA) ; ajustements de stock de FOUDA à Douala ; écart de
% 186 M FCFA ; registre R1 à R8 ; constats de l'audit des quatre centres.
% Faits nouveaux (reportés dans la bible) : demande du client PHARMACAM
% (fictif), étude EBIOS RM de janvier à mars 2027 animée par Mme NKENG ;
% ER1 à ER4 ; SR1 à SR5 ; parties prenantes P1 à P5 et leurs cotes ; base
% de preuve des 20 cotes de l'écosystème : 13 OUI / 1 NON / 6 NEK.
% ============================================================================

\chapter{EBIOS Risk Manager}
\label{chap:ebios}
\chaptermark{EBIOS Risk Manager}

\epigraph{%
    « On ne se protège pas d'une menace\\
    que l'on n'a pas pris la peine d'imaginer. »%
}{Principe du cours}

\niveauChapitre{Intermédiaire}{%
    Maîtrise du Chapitre~\ref{chap:iso27005}%
}

\begin{boxobjectifs}
\begin{itemize}
    \item Situer EBIOS Risk Manager~: son origine, sa logique « conformité
          et scénarios », ses cinq ateliers.
    \item Cadrer une étude~: valeurs métier, biens supports, événements
          redoutés, gravité, socle de sécurité.
    \item Identifier et retenir des couples source de risque et objectif
          visé.
    \item Cartographier l'écosystème et construire des scénarios
          stratégiques.
    \item Détailler un scénario opérationnel et en coter la vraisemblance.
    \item Coter la base de preuve d'une étude EBIOS, et en tirer les
          priorités de connaissance.
    \item Situer MEHARI et le NIST SP~800-30 par rapport à EBIOS.
\end{itemize}
\end{boxobjectifs}

\bigskip

% ============================================================================
\section*{Avant de commencer : le client demande une preuve}
% ============================================================================

Douala, janvier 2027. PHARMACAM\footnote{LiZbiZ, PHARMACAM, leurs
dirigeants, les dates et les montants sont fictifs~; la méthode et les
textes cités sont exacts.}, distributeur de médicaments et premier client
de LiZbiZ, prépare le renouvellement de son contrat de transport. Sa
direction des achats a ajouté une clause~: le prestataire devra remettre
« une analyse des risques de cybersécurité conduite selon une méthode
reconnue ». Les médicaments voyagent sous chaîne du froid~; un arrêt du
suivi des expéditions de plusieurs jours, et ce sont des lots entiers
qu'il faut détruire.

M.~ETEME, directeur général de LiZbiZ, appelle d'abord M.~BEYALA, d'AUDIX
CONSEIL, qui connaît l'entreprise. Celui-ci refuse d'animer l'étude~: son
cabinet a audité LiZbiZ et l'auditera encore~; il ne peut pas auditer
demain une analyse qu'il aurait conduite aujourd'hui. Il conseille de
confier l'animation à Mme~NKENG, contrôleuse interne, et de choisir EBIOS
Risk Manager.

\medskip

Ce chapitre présente cette méthode et la déroule sur LiZbiZ, atelier par
atelier. Les rançongiciels, la fraude au stock et les faux fournisseurs,
que le registre R1 à R8 décrivait séparément, y deviennent des scénarios
reliés à des attaquants, à leurs objectifs et à l'écosystème de
l'entreprise.

\begin{boxalert}
\textbf{Violation typique.} Faire animer l'analyse des risques par le
cabinet qui en auditera ensuite les résultats.

\smallskip
\textbf{Règle~:} on n'audite pas son propre travail. Un auditeur peut
conseiller sur la méthode~; il ne construit pas l'analyse qu'il devra
évaluer. Le Chapitre~\ref{chap:droit} revient sur ces règles
d'indépendance.
\end{boxalert}

\separateur

% ============================================================================
\section{La Méthode}
\label{sec:c04-methode}
% ============================================================================

\subsection{Origine et logique}

EBIOS Risk Manager est la méthode d'appréciation et de traitement du risque
numérique publiée en 2018 par l'Agence nationale de la sécurité des
systèmes d'information (\sigle{ANSSI}, France), avec le soutien du Club
EBIOS~\cite{anssi2018ebios}. Elle succède aux versions d'EBIOS de 1995,
2004 et 2010. Elle est compatible avec l'ISO~31000 et l'ISO/IEC~27005~:
elle en est une mise en œuvre, pas une alternative.

Sa logique tient en une phrase, que l'ANSSI formule comme une
\emph{synthèse entre conformité et scénarios}~:
\begin{itemize}
    \item la \textbf{conformité} construit un \emph{socle de sécurité}~:
          l'organisation applique les mesures d'un référentiel (hygiène
          informatique, ISO/IEC~27002) et mesure ses écarts~;
    \item les \textbf{scénarios} mettent ce socle à l'épreuve de menaces
          \emph{intentionnelles et ciblées}~: un attaquant précis, avec un
          objectif précis, qui passe éventuellement par l'écosystème de
          l'organisation.
\end{itemize}

\begin{boxinfo}
\textbf{Pourquoi cette division du travail~?} Les menaces courantes
(maliciel opportuniste, erreur, panne) se traitent par le socle~: il serait
inutile de construire un scénario pour chacune. L'effort d'imagination est
réservé aux attaques qui visent l'organisation et peuvent contourner le
socle. C'est l'inverse d'un registre qui énumère cent menaces au même
niveau de détail.
\end{boxinfo}

\subsection{Cinq ateliers}

\begin{tableaucours}{L{3.4cm}L{4.4cm}L{4.4cm}}{%
    \tableauHeader{Atelier} &
    \tableauHeader{Question} &
    \tableauHeader{Principaux livrables}
}
1. Cadrage et socle de sécurité & Que protège-t-on, de quoi a-t-on peur,
où en est-on de la conformité~? & Périmètre, valeurs métier, biens
supports, événements redoutés et gravité, écarts au socle \\
2. Sources de risque & Qui pourrait nous attaquer, et pour quoi faire~? &
Couples source de risque et objectif visé, retenus ou écartés \\
3. Scénarios stratégiques & Par où passerait l'attaquant, y compris par
nos partenaires~? & Cartographie de l'écosystème, scénarios stratégiques,
mesures sur l'écosystème \\
4. Scénarios opérationnels & Comment s'y prendrait-il, techniquement~? &
Modes opératoires, vraisemblance \\
5. Traitement du risque & Que décide-t-on, et comment le suit-on~? &
Stratégie de traitement, mesures, risques résiduels, cadre de suivi \\
\end{tableaucours}
\vspace{-0.8em}
\begin{center}{\small\itshape Tableau~4.1. Les cinq ateliers d'EBIOS Risk
Manager}\end{center}

Les ateliers~1 à~3 relèvent de la direction et des métiers~: ils
raisonnent en enjeux et en attaquants. L'atelier~4 relève des
techniciens~: il raisonne en vulnérabilités et en chemins d'attaque.
L'atelier~5 revient à la direction, qui décide. Les ateliers stratégiques se
revoient rarement, les ateliers opérationnels plus souvent, au rythme des
systèmes et des menaces.

\begin{boxinfo}
\textbf{Échelles.} Le guide propose des échelles de gravité et de
vraisemblance que chaque organisation adapte. Le cours emploie la sienne
(Chapitre~\ref{chap:risque}, Tableau~2.3), en appelant \emph{gravité}
l'impact d'un événement redouté, comme le fait la méthode. On note G1 à G4
et V1 à V4 les niveaux de cette échelle.
\end{boxinfo}

% ============================================================================
\section{Atelier 1 : Cadrage et Socle de Sécurité}
\label{sec:c04-atelier1}
% ============================================================================

\subsection{Périmètre, valeurs métier, biens supports}

Mme~NKENG réunit M.~ETEME, le directeur des opérations, le directeur
financier, M.~KOUNOU et un chef d'agence. Le périmètre retenu est le
système d'information logistique~: les quatre centres de données de Douala,
Yaoundé, Garoua et Bafoussam, les douze agences et l'application intégrée.
L'atelier retient trois valeurs métier~: le \emph{suivi des expéditions},
la \emph{facturation} et la \emph{gestion des stocks}, et une quatrième,
ajoutée à la demande de PHARMACAM~: les \emph{données des clients}.

\subsection{Événements redoutés et gravité}

\begin{tableaucours}{cL{5.6cm}L{4.2cm}c}{%
    \tableauHeader{Réf.} &
    \tableauHeader{Événement redouté} &
    \tableauHeader{Conséquences} &
    \tableauHeader{Grav.}
}
ER1 & Indisponibilité du suivi des expéditions pendant plus de trois jours
& Lots sous chaîne du froid perdus, pénalités, perte du client & G4 \\
ER2 & Altération des données de stock par des ajustements non autorisés &
Écarts de stock, réserve du commissaire aux comptes & G3 \\
ER3 & Paiement de factures à un fournisseur fictif & Perte financière
directe & G3 \\
ER4 & Divulgation des données des clients & Rupture de contrat,
sanction de l'autorité de protection & G3 \\
\end{tableaucours}
\vspace{-0.8em}
\begin{center}{\small\itshape Tableau~4.2. Événements redoutés de LiZbiZ}\end{center}

Les gravités s'appuient sur des faits~: l'écart de stock de 186~millions
FCFA (ER2) et la fraude TRANSLOG de 112,7~millions (ER3) se situent au
niveau majeur de l'échelle~; pour ER1, PHARMACAM a chiffré la valeur d'un
seul camion de vaccins à plus de 500~millions.

\subsection{Le socle de sécurité}

L'atelier compare ensuite l'existant à un référentiel. Faute de mieux,
LiZbiZ reprend les constats de l'audit de ses quatre centres~: des comptes
par défaut actifs, des services d'administration en clair, des flux non
autorisés entre centres. Ces écarts ne sont pas des scénarios~; ils
constituent un socle troué, que les scénarios des ateliers suivants vont
exploiter. Un écart au socle se corrige, en principe, sans attendre la fin
de l'étude.

% ============================================================================
\section{Atelier 2 : Sources de Risque}
\label{sec:c04-atelier2}
% ============================================================================

\subsection{Des couples, pas des menaces}

\begin{boxdef}
\termecle{Couple SR/OV}\index{couple SR/OV}~: association d'une source de
risque (qui~?) et d'un objectif visé (pour quoi faire~?). Un même
attaquant peut poursuivre plusieurs objectifs~; chaque couple se traite à
part.
\end{boxdef}

L'atelier recense les couples plausibles, puis les apprécie selon la
\emph{motivation} de la source à atteindre l'objectif, ses
\emph{ressources}, et son \emph{activité} connue dans le secteur. Il en
retient quelques-uns, jugés pertinents.

\begin{tableaucours}{cL{3.0cm}L{2.8cm}L{4.0cm}c}{%
    \tableauHeader{Réf.} &
    \tableauHeader{Source de risque} &
    \tableauHeader{Objectif visé} &
    \tableauHeader{Appréciation} &
    \tableauHeader{Retenu}
}
SR1 & Groupe cybercriminel & Extorsion par rançongiciel & Très motivé,
outillé, très actif contre les entreprises de transport & Oui \\
SR2 & Employé malveillant, seul ou avec un complice & Détourner du stock
et masquer l'écart & Précédent connu (ajustements de Douala) & Oui \\
SR3 & Fournisseur indélicat ou faux fournisseur & Se faire payer des
factures fictives & Précédent connu (TRANSLOG) & Oui \\
SR4 & Concurrent & Obtenir les tarifs et la liste des clients & Motivé,
mais peu d'activité connue & Non \\
SR5 & Militant & Nuire à l'image & Peu motivé par une entreprise de
logistique & Non \\
\end{tableaucours}
\vspace{-0.8em}
\begin{center}{\small\itshape Tableau~4.3. Couples source de risque et
objectif visé}\end{center}

\begin{boxscenario}{Opération LIZBIZ. Deux fraudes devenues des sources de risque}
\textbf{Constatation.} Les couples SR2 et SR3 ne sont pas imaginés~: ils
reprennent les deux fraudes établies par AUDIX CONSEIL, les ajustements de
stock de FOUDA et le fournisseur fictif TRANSLOG (112,7~millions FCFA).

\medskip
\textbf{Analyse} (confiance élevée). Le registre de 2026 traitait ces
fraudes comme des risques de contrôle interne (R1, R2), et le rançongiciel
comme un risque informatique (R8). EBIOS les met sur le même plan~: trois
attaquants, trois objectifs, qui exploitent souvent les mêmes faiblesses.
Les droits illimités de M.~KOUNOU servent aussi bien un fraudeur interne
qu'un cybercriminel qui aurait volé son compte.
\end{boxscenario}

% ============================================================================
\section{Atelier 3 : Scénarios Stratégiques}
\label{sec:c04-atelier3}
% ============================================================================

\subsection{L'écosystème}

Un attaquant ciblé passe souvent par un partenaire, moins protégé que sa
cible. L'atelier~3 cartographie donc les \emph{parties prenantes}~:
prestataires, sous-traitants, fournisseurs, clients reliés au système. Le
guide de l'ANSSI les évalue selon quatre critères~: la \emph{dépendance}
de l'organisation envers la partie prenante, la \emph{pénétration} de
celle-ci dans le système, sa \emph{maturité} en cybersécurité, et la
\emph{confiance} qu'on peut lui accorder. Il en déduit un niveau de
menace, d'autant plus élevé que la partie prenante est indispensable et
présente dans le système, et d'autant plus faible qu'elle est mûre et
fiable~:
\[
\text{Menace} = \frac{\text{Dépendance} \times \text{Pénétration}}
                     {\text{Maturité} \times \text{Confiance}}
\]
Chaque critère se cote de 1 à 4. L'atelier fixe lui-même les seuils~: une
menace de 3 ou plus place la partie prenante en zone de \emph{danger}, de
1 à 3 en zone de \emph{contrôle}, sous 1 en zone de \emph{veille}.

\begin{tableaucours}{cL{4.2cm}ccccc}{%
    \tableauHeader{Réf.} &
    \tableauHeader{Partie prenante} &
    \tableauHeader{Dép.} &
    \tableauHeader{Pén.} &
    \tableauHeader{Mat.} &
    \tableauHeader{Conf.} &
    \tableauHeader{Menace}
}
P1 & Éditeur de l'application intégrée (maintenance à distance) & 4 & 4 & 2
& 2 & 4,00 \\
P2 & Transporteurs sous-traitants (portail de suivi) & 3 & 2 & 1 & 2 &
3,00 \\
P3 & Opérateur de télécommunications & 4 & 1 & 3 & 3 & 0,44 \\
P4 & Cabinet comptable externe & 2 & 2 & 2 & 3 & 0,67 \\
P5 & Fournisseurs (portail de dépôt des factures) & 2 & 1 & 1 & 1 & 2,00 \\
\end{tableaucours}
\vspace{-0.8em}
\begin{center}{\small\itshape Tableau~4.4. Les parties prenantes de
LiZbiZ}\end{center}

L'éditeur de l'application (P1) et les transporteurs (P2) sont en zone de
danger~; les fournisseurs (P5) en zone de contrôle~; l'opérateur et le
cabinet comptable en zone de veille.

\subsection{Trois scénarios stratégiques}

\begin{boxdef}
\termecle{Scénario stratégique}\index{scénario stratégique}~: chemin
d'attaque de haut niveau, qui relie un couple SR/OV à un événement redouté,
en passant éventuellement par une ou plusieurs parties prenantes. Il se
décrit sans détail technique et prend la gravité de l'événement redouté.
\end{boxdef}

\begin{tableaucours}{cL{8.4cm}cc}{%
    \tableauHeader{Réf.} &
    \tableauHeader{Chemin d'attaque} &
    \tableauHeader{ER} &
    \tableauHeader{Grav.}
}
SS1 & SR1 compromet l'accès de maintenance de l'éditeur (P1), atteint les
centres de données et chiffre l'application et ses sauvegardes & ER1 & G4 \\
SS2 & SR2, magasinier de Douala, obtient l'aide d'un détenteur de droits
d'administration et passe des ajustements de stock sans contrôle & ER2 &
G3 \\
SS3 & SR3 s'enregistre comme fournisseur sur le portail (P5) et fait
valider des factures fictives par un complice interne & ER3 & G3 \\
\end{tableaucours}
\vspace{-0.8em}
\begin{center}{\small\itshape Tableau~4.5. Scénarios stratégiques de
LiZbiZ}\end{center}

L'atelier propose aussi des mesures sur l'écosystème, qui ne dépendent pas
du détail technique~: imposer à l'éditeur un second facteur et une
journalisation de ses connexions de maintenance~; limiter l'accès des
transporteurs aux seules expéditions qui les concernent~; vérifier
l'existence de tout nouveau fournisseur avant son premier paiement.

\subsection{Ce que valent les cotes de l'écosystème}

Les vingt cotes du Tableau~4.4 n'ont pas toutes la même solidité. La
dépendance et la pénétration se lisent dans les contrats et dans les
droits d'accès~; la maturité et la confiance se devinent.

\begin{boxlecture}{Base de preuve des cotes de l'écosystème}
\textbf{Constatation.} Mme~NKENG cote la base de chacune des vingt
cotes~:
\begin{itemize}
    \item dépendance~: 5~\OUI{} (contrats, chiffre d'affaires
          concerné)~;
    \item pénétration~: 4~\OUI{} et 1~\NON{}~: P2 était coté 2, mais les
          droits du portail montrent que chaque transporteur voit toutes
          les expéditions~;
    \item maturité~: 1~\OUI{} (l'opérateur est certifié ISO/IEC~27001), 4~\NEK{}~;
    \item confiance~: 3~\OUI{} (historique de la relation), 2~\NEK{} (P1, P2).
\end{itemize}
Sur vingt cotes~: 13~\OUI{}, 1~\NON{}, 6~\NEK{}~:
\[
\mu = 0{,}65~;\quad U = 0{,}30~;\quad
H = -(0{,}65\log_2 0{,}65 + 0{,}05\log_2 0{,}05 + 0{,}30\log_2 0{,}30)
  \approx 0{,}404 + 0{,}216 + 0{,}521 \approx 1{,}14~\text{bit}.
\]

\textbf{Analyse} (confiance élevée sur le calcul, moyenne sur la lecture).
L'incertitude se concentre au dénominateur de la formule~: la maturité et
la confiance. Or ce sont elles qui font baisser la menace. Une maturité
cotée 2 sans preuve peut cacher un 1, ce qui doublerait la menace de P1. Les
zones de danger sont donc probablement sous-estimées, et non surestimées.
Pour P2, le \NON{} se corrige tout de suite~: avec une pénétration de 3, sa
menace passe de 3,00 à 4,50.

\medskip
\hfill\lecture{0{,}65}{0{,}30}{1{,}14}
\end{boxlecture}

\begin{boxprocedure}
\textbf{Réduire l'incertitude sur l'écosystème.} Pour chaque partie
prenante en zone de danger ou de contrôle, le papier de travail prévoit~:
\begin{itemize}
    \item un questionnaire de maturité, signé par la partie prenante, et
          les attestations qu'elle peut fournir (certification, rapport
          d'assurance)~;
    \item l'extraction des droits d'accès dont elle dispose réellement~;
    \item une clause contractuelle de droit d'audit et de notification des
          incidents~;
    \item la date de la dernière revue, et la cote de base de chaque
          critère.
\end{itemize}
\end{boxprocedure}

% ============================================================================
\section{Atelier 4 : Scénarios Opérationnels}
\label{sec:c04-atelier4}
% ============================================================================

\subsection{Du chemin au mode opératoire}

\begin{boxdef}
\termecle{Scénario opérationnel}\index{scénario opérationnel}~: suite
d'actions élémentaires, sur les biens supports, par lesquelles une source
de risque réalise un scénario stratégique. On le décompose en quatre
phases~: \emph{connaître} la cible, y \emph{rentrer}, y \emph{trouver} ce
qui permet d'atteindre l'objectif, puis l'\emph{exploiter}.
\end{boxdef}

\subsection{Le mode opératoire de SS1}

M.~KOUNOU et deux techniciens de l'éditeur détaillent SS1. Pour chaque
phase, ils cotent la vraisemblance que l'attaquant y réussisse, au vu des
constats de l'audit.

\begin{tableaucours}{L{1.8cm}L{6.2cm}cL{2.6cm}}{%
    \tableauHeader{Phase} &
    \tableauHeader{Action élémentaire} &
    \tableauHeader{V} &
    \tableauHeader{Base}
}
Connaître & Repérer sur Internet l'éditeur et son outil de maintenance à
distance & V4 & Sources publiques \\
Rentrer & Voler l'identifiant d'un technicien de l'éditeur par
hameçonnage, ou utiliser un compte par défaut du serveur de maintenance &
V3 & Constat d'audit (comptes par défaut) \\
Trouver & Se déplacer d'un centre à l'autre par les flux non autorisés,
capter des mots de passe sur les services en clair & V4 & Constats
d'audit \\
Exploiter & Chiffrer l'application dans les quatre centres et effacer les
sauvegardes accessibles & V3 & Sauvegardes en ligne, jamais restaurées \\
\end{tableaucours}
\vspace{-0.8em}
\begin{center}{\small\itshape Tableau~4.6. Le scénario opérationnel SO1}\end{center}

\begin{boxinfo}
\textbf{Une règle de chaîne.} Un mode opératoire ne réussit que si toutes
ses phases réussissent. Le cours retient donc, pour vraisemblance du
scénario, celle de la phase la plus difficile pour l'attaquant~: ici V3.
C'est une convention prudente~; le guide décrit aussi des manières plus
fines de combiner les actions élémentaires.
\end{boxinfo}

SO1 obtient une vraisemblance V3 et hérite de la gravité G4 de ER1~:
criticité $3 \times 4 = 12$, intolérable. On retrouve la cote de R8 après
sa révision (Exercice~2.4)~: EBIOS ne la contredit pas, il l'explique. Et
il montre où agir~: la phase «~rentrer~» est la seule cotée V3 au lieu de
V4~; la fermer (second facteur pour l'éditeur, suppression des comptes par
défaut) ferait tomber la vraisemblance de tout le scénario.

% ============================================================================
\section{Atelier 5 : Traitement du Risque}
\label{sec:c04-atelier5}
% ============================================================================

L'atelier~5 rejoint le processus du Chapitre~\ref{chap:iso27005}~: options
de traitement, plan, risques résiduels, acceptation. EBIOS y ajoute deux
apports~: les mesures se rattachent aux \emph{phases} des scénarios et aux
\emph{parties prenantes}, et le plan s'inscrit dans un cadre de suivi qui
fait revenir l'étude à intervalles réguliers. Le guide parle d'un plan
d'amélioration continue de la sécurité.

\begin{tableaucours}{cL{5.6cm}L{2.4cm}c}{%
    \tableauHeader{Scénario} &
    \tableauHeader{Mesures principales} &
    \tableauHeader{Phase ou partie prenante} &
    \tableauHeader{Résiduel visé}
}
SO1 (SS1) & Second facteur et journalisation pour l'éditeur~; suppression
des comptes par défaut~; filtrage entre centres~; sauvegardes hors ligne
testées & Rentrer, trouver, exploiter~; P1 & $2 \times 3 = 6$ \\
SS2 & Séparation des tâches de réception et d'ajustement~; revue des
ajustements par la contrôleuse interne~; comptes d'administration
nominatifs & Exploiter & $2 \times 3 = 6$ \\
SS3 & Vérification de tout nouveau fournisseur avant paiement~; double
validation des factures & P5 & $1 \times 3 = 3$ \\
\end{tableaucours}
\vspace{-0.8em}
\begin{center}{\small\itshape Tableau~4.7. Synthèse du traitement}\end{center}

\begin{boxobjet}{SO1 · Rançongiciel par l'accès de maintenance}
\textbf{Valeur métier~:} suivi des expéditions.

\textbf{Événement redouté~:} ER1, indisponibilité de plus de trois jours
(G4).

\textbf{Source de risque et objectif visé~:} SR1, groupe cybercriminel,
extorsion.

\textbf{Partie prenante~:} P1, éditeur de l'application (menace 4,00,
maturité et confiance \NEK{}).

\textbf{Mode opératoire~:} connaître (V4), rentrer (V3), trouver (V4),
exploiter (V3)~; vraisemblance retenue V3.

\textbf{Criticité~:} 12, intolérable. Base~: constats de l'audit des quatre
centres (\OUI{}) pour les phases rentrer, trouver et exploiter.

\textbf{Traitement~:} réduire (Tableau~4.7), puis envisager une assurance.

\textbf{Propriétaire~:} le directeur des opérations.

\textbf{Revue~:} après le test des sauvegardes et la signature de
l'avenant avec l'éditeur, et au plus tard dans un an.
\end{boxobjet}

% ============================================================================
\section{Situer EBIOS parmi les Méthodes}
\label{sec:c04-methodes}
% ============================================================================

EBIOS Risk Manager n'est pas la seule méthode. Deux autres sont souvent
citées~: MEHARI, développée par le CLUSIF, l'association française des
professionnels de la sécurité, et le guide d'appréciation des risques du
NIST, la publication spéciale 800-30~\cite{nist2012sp80030}.

\begin{tableaucours}{L{2.4cm}L{3.3cm}L{3.3cm}L{3.3cm}}{%
    \tableauHeader{} &
    \tableauHeader{EBIOS RM} &
    \tableauHeader{MEHARI} &
    \tableauHeader{NIST SP~800-30}
}
Origine & ANSSI, 2018 & CLUSIF & NIST, révision~1 de 2012 \\
Point de départ & Sources de risque, objectifs, écosystème & Bases de
connaissances, diagnostic des services de sécurité & Sources de menace,
événements de menace, vulnérabilités \\
Force & Scénarios lisibles par la direction~; écosystème & Questionnaires
détaillés, utiles à l'audit & Démarche à trois niveaux (organisation,
mission, système) \\
Limite & Demande des animateurs formés & Lourde sans outil & Pensée pour
l'administration américaine~; à adapter \\
\end{tableaucours}
\vspace{-0.8em}
\begin{center}{\small\itshape Tableau~4.8. Trois méthodes d'appréciation
des risques}\end{center}

Pour une organisation camerounaise, le choix n'est pas dicté par la loi~:
l'article~26 de la loi n°~2010/012 exige des « systèmes normalisés »
d'identification et d'évaluation des risques (Chapitre~\ref{chap:iso27005}),
sans désigner de méthode~\cite{loi2010012}. EBIOS a l'avantage d'être
gratuite, publiée en français et compatible avec l'ISO/IEC~27005~; elle
reste une méthode étrangère, que l'on adapte aux échelles, au droit et aux
menaces du pays.

% ============================================================================
\section{Ce que la Méthode demande à chaque Métier}
\label{sec:c04-metiers}
% ============================================================================

\begin{boxmetier}{Gestionnaire des risques cyber}
\textbf{Animer, pas rédiger seul}
\begin{enumerate}
    \item Réunir les bons participants à chaque atelier~: la direction et
          les métiers pour les ateliers~1 à~3, les techniciens pour le~4.
    \item Partir des faits connus (fraudes, incidents, constats d'audit)
          pour les couples SR/OV.
    \item Coter la base de preuve des critères de l'écosystème, et
          programmer les questionnaires et les extractions qui manquent.
    \item Relier chaque mesure à une phase ou à une partie prenante.
\end{enumerate}
\end{boxmetier}

\begin{boxmetier}{Auditeur en cybersécurité}
\textbf{Rester de l'autre côté}

M.~BEYALA a eu raison de refuser l'animation. L'auditeur peut évaluer une
étude EBIOS~: il vérifie que les couples écartés l'ont été pour de bonnes
raisons, que les cotes de l'écosystème ont une base, et que les mesures
annoncées existent. Il ne l'écrit pas.
\end{boxmetier}

\begin{boxmetier}{Responsable GRC}
\textbf{Ramener l'étude dans le registre}

Une étude EBIOS produit des scénarios~; le registre doit les recevoir. Le
responsable GRC remplace les lignes R1, R2 et R8 par les scénarios SS2,
SS3 et SO1, avec leurs propriétaires, et veille à ce que l'étude soit
revue quand un nouveau partenaire entre dans l'écosystème.
\end{boxmetier}

\begin{boxmetier}{Consultant en mise en œuvre de SMSI}
\textbf{Transmettre la méthode}

Quand un consultant anime l'étude, il forme en même temps un animateur
interne. La demande de PHARMACAM se répétera chaque année~; LiZbiZ doit
pouvoir refaire l'étude sans lui.
\end{boxmetier}

\separateur

\begin{boxresume}
\textbf{Ce qu'il faut retenir de ce chapitre~:}
\begin{itemize}
    \item \textbf{EBIOS Risk Manager} (ANSSI, 2018) réalise une synthèse
          entre conformité (un socle de sécurité) et scénarios (des menaces
          intentionnelles et ciblées).
    \item \textbf{Cinq ateliers}~: cadrage et socle, sources de risque,
          scénarios stratégiques, scénarios opérationnels, traitement.
    \item \textbf{Un couple SR/OV} relie un attaquant à un objectif~; on
          part des faits connus pour les retenir.
    \item \textbf{L'écosystème} se cote par la dépendance, la pénétration,
          la maturité et la confiance~; les deux derniers critères sont
          souvent \NEK{}, et ce sont eux qui font baisser la menace.
    \item \textbf{Un scénario opérationnel} se décompose en connaître,
          rentrer, trouver, exploiter~; sa vraisemblance est celle de la
          phase la plus difficile.
    \item \textbf{Indépendance}~: celui qui auditera l'étude ne l'anime pas.
\end{itemize}
\end{boxresume}

\begin{boxinfo}
\textbf{Points de contrôle.} Vous avez compris ce chapitre si vous savez~:
\begin{itemize}
    \item dire à quel atelier appartient un livrable donné~;
    \item calculer le niveau de menace d'une partie prenante et la placer
          dans une zone~;
    \item décomposer un scénario opérationnel en quatre phases et en coter
          la vraisemblance~;
    \item expliquer pourquoi une maturité cotée sans preuve sous-estime
          probablement la menace.
\end{itemize}
\end{boxinfo}

% ============================================================================
\section*{Exercices}
% ============================================================================

\begin{boxexo}[title={\ding{45}\quad Exercice~4.1. À quel atelier~? \niveaubase}]
Indiquez l'atelier d'EBIOS Risk Manager auquel appartient chacun des
éléments suivants.
\begin{enumerate}[(a)]
    \item La liste des écarts de LiZbiZ au guide d'hygiène informatique.
    \item Le couple « concurrent, obtenir la liste des clients », écarté.
    \item La cote de pénétration du cabinet comptable.
    \item L'action « capter des mots de passe sur les services en clair ».
    \item La gravité G4 de l'indisponibilité du suivi des expéditions.
    \item La décision de faire signer un avenant à l'éditeur.
    \item Le chemin « faux fournisseur, portail, complice interne ».
\end{enumerate}
\end{boxexo}

\begin{boxexo}[title={\ding{45}\quad Exercice~4.2. Trois nouvelles parties prenantes \niveaumoyen}]
LiZbiZ envisage trois nouveaux partenaires~:
\begin{itemize}
    \item Q1, un hébergeur en nuage pour le portail de suivi~:
          dépendance~4, pénétration~3, maturité~3, confiance~2~;
    \item Q2, une société de maintenance des chambres froides, reliée au
          réseau pour la télésurveillance~: dépendance~3, pénétration~2,
          maturité~1, confiance~2~;
    \item Q3, un prestataire de paie~: dépendance~2, pénétration~1,
          maturité~2, confiance~3.
\end{itemize}
\begin{enumerate}[(a)]
    \item Calculez leur niveau de menace et placez-les dans les zones du
          chapitre.
    \item Les maturités ont été cotées sur la base d'une plaquette
          commerciale. Quelle cote de base leur donnez-vous~? Recalculez la
          menace de Q1 si sa maturité réelle est de~2.
    \item Proposez, pour Q2, deux mesures sur l'écosystème et la preuve
          qui permettrait de relever sa confiance.
\end{enumerate}
\end{boxexo}

\begin{boxexo}[title={\ding{45}\quad Exercice~4.3. Le mode opératoire de SS2 \niveaumoyen}]
Détaillez SS2 (le magasinier de Douala et son complice) en scénario
opérationnel.
\begin{enumerate}[(a)]
    \item Décrivez une action élémentaire pour chacune des quatre phases.
    \item Cotez la vraisemblance de chaque phase et justifiez par un fait du
          cas. Déduisez la vraisemblance du scénario selon la règle de
          chaîne.
    \item Quelle phase faut-il fermer en priorité~? Montrez l'effet sur la
          criticité.
\end{enumerate}
\end{boxexo}

\begin{boxexo}[title={\ding{45}\quad Exercice~4.4. EBIOS à la BANQUE SANAGA \niveauexpert}]
L'atelier de décembre 2026 de la BANQUE SANAGA a retenu le scénario
« fraude au virement par compromission d'un poste du siège »
(Chapitre~\ref{chap:iso27005}).
\begin{enumerate}[(a)]
    \item Atelier~1~: nommez la valeur métier, deux biens supports et
          l'événement redouté, et cotez sa gravité avec l'échelle du cours.
    \item Atelier~2~: proposez deux couples SR/OV, dont un que vous
          écarterez, en justifiant par la motivation, les ressources et
          l'activité.
    \item Atelier~3~: identifiez une partie prenante qui pourrait servir de
          point d'entrée, cotez ses quatre critères, et indiquez pour chacun
          la base de preuve dont vous disposez.
    \item Atelier~4~: rédigez le mode opératoire en quatre phases et
          déduisez la vraisemblance.
    \item Comparez la démarche avec la fiche R-AG-01~: qu'apporte EBIOS que
          la fiche de risque ne donnait pas~?
\end{enumerate}
\end{boxexo}

\bigskip

\begin{boxinfo}
\textbf{Pour aller plus loin.}
\begin{itemize}
    \item Le guide EBIOS Risk Manager de l'ANSSI~\cite{anssi2018ebios} et ses
          fiches méthodes, sur le site de l'ANSSI.
    \item Le guide NIST SP~800-30, révision~1~\cite{nist2012sp80030}.
    \item L'ISO/IEC~27005:2022~\cite{iso27005_2022}, pour l'articulation
          avec le processus.
\end{itemize}
\end{boxinfo}

\begin{boxinfo}
\textbf{Transition.} EBIOS a coté SO1 à 12~: intolérable. Mais combien
LiZbiZ ou la BANQUE SANAGA perdraient-elles, en francs CFA, si le scénario
se réalisait~? Et combien faut-il dépenser pour le traiter~? Les échelles
qualitatives ne répondent pas à ces questions. Le Chapitre~\ref{chap:fair}
apprend à quantifier.
\end{boxinfo}
